\documentclass[UTF8,fontset=none,11pt,a4paper]{ctexart}
\usepackage{ipara-handbook}

\handbooksetup
  {局部到整体：中值定理与函数形状}
  {数学一 · 高等数学 Topic 04}
  {中值定理 · 零点 · 函数形状}

\begin{document}
\handbooktitle
  {局部到整体：中值定理与函数形状}
  {从“存在一个中值点”到“恢复整个区间的行为”}
  {}

\begin{quote}
本册只追问一个根本问题：\textbf{已知函数在端点、零点或局部导数上的信息，怎样把它运输成区间内部的见证点和整体形状？} 中值定理负责产生内部证据，辅助函数负责把目标表达式翻译成可用的导数结构，导数符号再把点态信息恢复成单调、极值、凹凸和零点分布。
\end{quote}

\begin{center}\rule{0.5\linewidth}{0.5pt}\end{center}

% ============================================================
\setcounter{section}{-1}
\section{本册定位}
% ============================================================

本册是高等数学 Subject Atlas 下的 Topic04。Topic02 提供连续与极限资格，Topic03 提供导数和有限局部模型；本册研究这些局部/边界信息怎样借助区间定理产生整体结论。

\subsection{Own}

本册独立负责：
\begin{itemize}
  \item 介值定理、零点定理、Rolle、Lagrange 与 Cauchy 中值定理的资格和逻辑链；
  \item 证明题中由目标表达式反推辅助函数的构造机制；
  \item 零点、重零点与反复 Rolle 的区间见证；
  \item 导数符号到单调性、极值、凹凸、拐点和函数形状的恢复；
  \item 驻点、临界点、极值点和拐点的候选/确认边界；
  \item 用零点重数、导数根的符号与 Rolle 链分析根的数量；
  \item 中值定理证明与函数形状题的统一检查协议。
\end{itemize}

\subsection{Uses / Stop Boundary}

本册调用 Topic01 的函数对象/定义域，Topic02 的连续性，Topic03 的局部导数和有限 Taylor。以下内容只保留接口：
\begin{itemize}
  \item 多中值点插点法、中值点参数渐近位置、插值余项与导数估计归 H-B02；
  \item Taylor 多项式本体归 Topic03；Lagrange 型余项的区间误差接口由 H-B02 接收；
  \item 变上限积分、原函数和积分因子所需的正则性归 Topic05/H-B03；
  \item 一阶微分方程的整体解族归 Topic12，本册只把积分因子作为辅助函数构造的局部工具；
  \item 题目看到什么信号、如何排序和何时退出进入《高等数学做题规则》。
\end{itemize}

% ============================================================
\section{根本问题：端点信息怎样变成区间内部证据？}
% ============================================================

只知道 $f(a)$、$f(b)$ 或 $f'(a)$，通常不能直接写出某个内部点 $\xi$。朴素做法是“找一个好看的点”或直接把平均斜率当作某点斜率，但这缺少存在性证明，也没有解释为什么目标表达式会在该点出现。

中值定理提供一条受资格条件保护的运输通道：
\begin{itemize}
  \item 连续性把端点值之间的中间输出全部填满；
  \item 端点等值把函数压成一条必须在内部水平的轨迹；
  \item 端点不等时，减去割线后制造端点等值，得到平均变化率；
  \item 两个函数同时变化时，用 Cauchy 的共享参数运输比值关系；
  \item 足够多个零点或重零点经过反复 Rolle，产生高阶导数见证。
\end{itemize}

因此区间定理不是“记住四个定理”，而是把\textbf{边界约束转译成内部见证}。

\begin{corebox}{中值点只保证存在，不给出可随意选择的固定坐标}
结论“存在 $\xi\in(a,b)$”是一个存在性输出。除非题设进一步唯一化或给出估计，不能把 $\xi$ 当成某个端点、对称点或自己挑选的数。所有回译都必须保留 $\xi$ 的合法区间和分母条件。
\end{corebox}

% ============================================================
\section{母模型：目标、辅助函数、资格、见证与回译}
% ============================================================

本册把中值定理和函数形状压缩为六个依次发生的动作：
\begin{processchain}
  \processstage{Target}
  \processstage{Auxiliary Function}
  \processstage{Qualification}
  \processstage{Boundary/Zero Design}
  \processstage{Witness}
  \processstage{Translate Back}
\end{processchain}

这里的箭头表示\textbf{证明和分析的控制顺序}：
\begin{enumerate}
  \item \kw{Target}：把待证结论写成 $\Phi(\xi)=0$ 或某个导数/符号目标；
  \item \kw{Auxiliary Function}：反推一个 $F$，使 $F^{(k)}$ 等于 $\Phi$ 的非零倍数；
  \item \kw{Qualification}：核对闭区间连续、开区间可导、分母非零、对数真数和积分因子；
  \item \kw{Boundary/Zero Design}：让 $F$ 端点等值，或制造足够多个不同零点/重零点；
  \item \kw{Witness}：应用介值、Rolle、Lagrange、Cauchy 或反复 Rolle，产生合法内部点；
  \item \kw{Translate Back}：展开 $F^{(k)}(\xi)=0$，用非零因子回译为原目标，并保留区间信息。
\end{enumerate}

\begin{center}
\begin{tikzpicture}[
  node distance=0.45cm and 0.48cm,
  box/.style={draw=iparaDeep,rounded corners=1.5mm,minimum width=3.0cm,minimum height=1.0cm,align=center,fill=iparaSoft,font=\small},
  arr/.style={-{Latex[length=2mm]},thick,draw=iparaDeep}
]
\node[box] (a) {目标表达式\\$\Phi(\xi)=0$};
\node[box,right=of a] (b) {反推辅助函数\\$F^{(k)}=m\Phi$};
\node[box,right=of b] (c) {资格审查\\连续/可导/非零};
\node[box,below=of c] (d) {制造约束\\端点等值/多零点};
\node[box,left=of d] (e) {中值见证\\$\xi$ 存在};
\node[box,left=of e] (f) {回译结论\\保留区间与条件};
\draw[arr] (a)--(b);
\draw[arr] (b)--(c);
\draw[arr] (c)--(d);
\draw[arr] (d)--(e);
\draw[arr] (e)--(f);
% Fallback 是外围反馈，不与主流程争用底部中心入射空间。
% a.south 下方已有 f，故 south 属于 blocked side；改从 a.west 的 free side 接入。
\coordinate (fallbackBottom) at ($(d.south east)+(0,-0.42)$);
\coordinate (fallbackLeft) at ($(a.west)+(-0.42,0)$);
\draw[arr,dashed,draw=iparaWarn]
  (d.south east) -- (fallbackBottom)
  -- node[midway,below,font=\scriptsize]{约束不足：停止套定理} (fallbackBottom -| fallbackLeft)
  -- (fallbackLeft) -- (a.west);
\end{tikzpicture}
\end{center}

函数形状题运行同一模型的逆向版本：先由 $f'$ 或 $f''$ 的符号得到局部/区间见证，再通过 Lagrange 把符号运输成函数值差异，最后确认极值或凹凸边界。

% ============================================================
\section{定理阶梯：连续、水平、割线与双函数运输}
% ============================================================

\subsection{介值定理与零点定理：连续性填补输出空隙}

若 $f$ 在 $[a,b]$ 上连续，则任意介于 $f(a)$ 与 $f(b)$ 的值 $\mu$ 都由某个 $c\in[a,b]$ 取得。若 $f(a)f(b)<0$，则存在 $c\in(a,b)$ 使 $f(c)=0$。

这不是“图像看起来会穿过横轴”的直觉，而是连续性阻止跳跃的存在性结论。若端点值恰好为零，见证点可能在端点；要求内部零点时必须明确端点非零或改变区间。

\subsection{Rolle：端点等值制造内部水平切线}

若 $f$ 在 $[a,b]$ 连续、在 $(a,b)$ 可导且 $f(a)=f(b)$，则存在 $\xi\in(a,b)$ 使
\[
f'(\xi)=0.
\]

Rolle 是整个定理链的发动机：只要能把一个目标改写成“某个辅助函数两端相等”，就能获得内部导数为零的见证。端点等值是必须验证的边界约束，不是可以省略的形式。

\subsection{Lagrange：减去割线后恢复平均变化率}

若 $f$ 在 $[a,b]$ 连续、在 $(a,b)$ 可导，则取割线斜率
\[
K=\frac{f(b)-f(a)}{b-a},
\]
构造
\[
F(x)=f(x)-f(a)-K(x-a).
\]
此时 $F(a)=F(b)=0$。由 Rolle 得到
\[
F'(\xi)=f'(\xi)-K=0,
\]
即
\[
\boxed{f'(\xi)=\frac{f(b)-f(a)}{b-a}}.
\]

Lagrange 的机制不是“端点斜率等于某点斜率”，而是先减掉已知的线性趋势，把问题变成水平端点问题。

\subsection{Cauchy：共同参数下的比值运输}

若 $f,g$ 在 $[a,b]$ 连续、在 $(a,b)$ 可导，则存在 $\xi$ 使
\[
[f(b)-f(a)]g'(\xi)=[g(b)-g(a)]f'(\xi).
\]
只有在进一步确保 $g'(\xi)\ne0$ 或对应差值非零时，才可把它改写为导数比。安全证明应优先保留交叉相乘形式，避免未经资格审查的除法。

可构造
\[
F(x)=[f(b)-f(a)][g(x)-g(a)]-[g(b)-g(a)][f(x)-f(a)],
\]
其两端均为零，再应用 Rolle。Cauchy 是 Lagrange 的双函数接口，不是“看到分式就自动洛必达”。

\subsection{高阶 Rolle：零点数量转为高阶导数见证}

若 $F$ 有 $n+1$ 个互不相同的零点，并具备相应阶数的连续/可导条件，则反复应用 Rolle 可得某个 $\xi$ 使
\[
F^{(n)}(\xi)=0.
\]
重零点按其导数消失阶数计数。例如 $F(a)=F'(a)=0$ 等价于在反复 Rolle 中提供两个零点条件。中间每一次得到的点都位于相邻零点之间，不能把不同层的中值点混为同一个点。

% ============================================================
\section{证明题逆向工程：从目标表达式反推辅助函数}
% ============================================================

\subsection{先把结论标准化}

把“存在 $\xi$ 使某式成立”写成
\[
\Phi(\xi)=0.
\]
然后问：能否找到 $F$ 与在区间内不为零的 $m$，使
\[
F'(x)=m(x)\Phi(x)
\]
或更高阶地 $F^{(k)}(x)=m(x)\Phi(x)$？这一步把猜辅助函数变成了可检查的逆向设计。

\subsection{结构反推表}

\begin{handbooklongtable}{L{0.20\linewidth}L{0.28\linewidth}Y}
\toprule
\textbf{目标导数结构} & \textbf{辅助函数} & \textbf{必须核对的资格}\\
\midrule
$f'(x)=K$ & $f(x)-Kx$ 或减去割线 & 端点等值、区间连续/可导\\
$A'B+AB'$ & $AB$ & 因子在区间内有定义\\
$A'B-AB'$ & $A/B$ & $B$ 在区间内不为零\\
$A'/A\pm B'/B$ & $\ln|A|\pm\ln|B|$ & $A,B$ 不为零且符号区间固定\\
$y'+py=q$ & 积分因子乘积减变限积分 & $p,q$ 连续，指数与积分合法\\
$F^{(n)}(x)=K$ & 减去满足插值条件的低次多项式 & 足够多个零点/重零点\\
\bottomrule
\end{handbooklongtable}

表格提供的是机制接口，不是无条件题型口诀。每次使用仍须先完成 Qualification 和 Boundary\allowbreak/Zero Design 两步。

\subsection{乘积、商与对数结构}

若目标含 $A'B+AB'$，逆用乘积法则取 $F=AB$；若含 $A'B-AB'$ 且 $B\ne0$，取 $F=A/B$。若含 $A'/A$，对数辅助函数可把乘法结构压成加法，但真数非零与符号保持是不可删除的资格。

例如要证明存在 $\xi$ 使
\[
f'(\xi)[g(\xi)-g(b)]+g'(\xi)[f(\xi)-f(a)]=0,
\]
取
\[
F(x)=[f(x)-f(a)][g(x)-g(b)].
\]
有 $F(a)=F(b)=0$，Rolle 给出 $F'(\xi)=0$，展开即得目标。真正的设计动作是让两个因子分别在相反端点消失。

\subsection{积分因子与变上限函数只作局部工具}

对一阶线性结构 $y'+p(x)y=q(x)$，令
\[
\mu(x)=\exp\left(\int p(x)\,dx\right)
\]
使 $[\mu y]'=\mu(y'+py)$。非齐次目标可以再减去一个变上限积分，使辅助函数的导数等于方程残差的非零倍数。

但“总能定义 $F(x)=\int\Phi$”不自动产生中值点；还必须另外证明 $F$ 在两个点等值或拥有足够零点。积分只解决导数匹配，不替代边界设计。

% ============================================================
\section{函数形状：把导数符号运输成区间性质}
% ============================================================

\subsection{单调性是 Lagrange 的全区间回译}

若 $f$ 在区间 $I$ 上连续、在内部可导，且对任意 $x$ 有 $f'(x)>0$，任取 $x_1<x_2$，由 Lagrange 存在 $\xi\in(x_1,x_2)$ 使
\[
f(x_2)-f(x_1)=f'(\xi)(x_2-x_1)>0.
\]
所以 $f$ 严格递增；$f'<0$ 同理严格递减；$f'=0$ 则函数为常数。

导数符号是充分条件。严格递增函数在某些点可能导数为零，不能把“递增”反向升级为“处处导数严格为正”。

\subsection{极值：先列候选，再检查符号变化}

内部极值点若可导，必满足 $f'(a)=0$；不可导尖点也可能是极值。因此候选集至少包括：
\begin{itemize}
  \item 驻点 $f'(a)=0$；
  \item $f'$ 不存在但函数有定义的点；
  \item 定义域端点（若问闭区间绝对极值）。
\end{itemize}

候选不等于结论。最稳的确认方法是检查 $f'$ 在两侧的符号：负到正为极小，正到负为极大，符号不变则不是极值。若 $f'(a)=0$ 且 $f''(a)>0$ 或 $<0$，二阶充分条件可快速确认；$f''(a)=0$ 时必须回到符号变化或更高阶第一非零导数。

\subsection{凹凸与拐点：二阶符号也要看区间}

在二阶可导区间内，$f''>0$ 表示斜率递增、图像凹向上；$f''<0$ 表示斜率递减、图像凹向下。拐点的核心是凹凸性在点两侧发生改变，不是单独看到 $f''(a)=0$。

若 $f''(a)=0$ 且 $f^{(3)}(a)\ne0$，在相应连续性资格下通常可确认拐点；若二阶导数在点不存在，仍可通过两侧凹凸定义判断。不可导点可以是拐点，但必须检查函数连续和凹凸真正改变。

\subsection{三点割线序与导函数单调：严格性要分两层恢复}

二阶导数只是判断凸性的一条充分路径，不是唯一入口。若 $f$ 在 $(a,b)$ 可导，则下面两件事等价：
\begin{enumerate}
  \item $f'$ 在 $(a,b)$ 严格递增；
  \item 任意 $x_1<x_2<x_3$ 都满足相邻割线斜率严格递增：
  \[
  \frac{f(x_2)-f(x_1)}{x_2-x_1}
  <
  \frac{f(x_3)-f(x_2)}{x_3-x_2}.
  \]
\end{enumerate}

由 1 推 2 很直接：分别在 $[x_1,x_2]$ 与 $[x_2,x_3]$ 使用 Lagrange，得到
$\xi_1\in(x_1,x_2)$、$\xi_2\in(x_2,x_3)$。因为 $\xi_1<\xi_2$，两条割线斜率分别等于
$f'(\xi_1)$ 与 $f'(\xi_2)$，严格次序随即成立。

反向不能把严格割线不等式直接``取极限''成严格导数不等式。任取 $u<v$，再取
$x<u<v<y$，由三点条件两次得到
\[
\frac{f(u)-f(x)}{u-x}
<
\frac{f(v)-f(u)}{v-u}
<
\frac{f(y)-f(v)}{y-v}.
\]
令 $x\to u^-$、$y\to v^+$，只能推出
\[
f'(u)\le
\frac{f(v)-f(u)}{v-u}
\le f'(v),
\]
所以第一层结论只是 $f'$ 非减。

严格性来自第二层的\textbf{平台排除}。若某个 $u<v$ 满足 $f'(u)=f'(v)$，非减性迫使
$f'$ 在 $[u,v]$ 上恒定；由 Lagrange，$f$ 在该段为仿射函数。于是在 $(u,v)$ 内任取三个点，相邻割线斜率相等，与严格三点条件矛盾。因此任意 $u<v$ 都有 $f'(u)<f'(v)$。

\begin{warnbox}{严格不等式经过极限可能降级}
若 $A_x<B<C_y$，即使 $A_x\to A$、$C_y\to C$，通常也只能得到 $A\le B\le C$。本题必须先得到导函数非减，再另用``若端点导数相等则整段为平台''恢复严格性。直接写``令端点趋近，所以 $f'(u)<f'(v)$''缺少关键证明。
\end{warnbox}

\subsection{零点重数与符号穿越}

若
\[
f(x)={(x-a)}^n g(x),
\qquad g(a)\ne0,
\]
则 $n$ 的奇偶决定函数在 $a$ 两侧是否变号：偶数阶通常保号并产生极值型接触，奇数阶变号；当奇数阶至少为三时，常伴随凹凸翻转。这里“通常/常伴随”不能替代直接检查 $g$ 的符号与凹凸。

对于多项式，$f'$ 的奇数重根会导致 $f'$ 变号并给出极值横坐标；$f''$ 的奇数重根会导致凹凸翻转并给出拐点横坐标。这个技巧依赖多项式/连续因子结构，不能无条件推广到带振荡或不规则因子的函数。

% ============================================================
\section{贯穿母例：三次函数的一条完整区间轨迹}
% ============================================================

令
\[
f(x)=x^3-3x.
\]
这是一个同时暴露候选点、符号运输、形状边界和辅助函数构造的母例。

\subsection{从导数符号到形状}

\[
f'(x)=3(x^2-1)=3(x-1)(x+1).
\]
因此 $f'>0$ 于 $(-\infty,-1)$ 与 $(1,\infty)$，$f'<0$ 于 $(-1,1)$。Lagrange 回译出：函数先增、后减、再增；$x=-1$ 是极大点，$x=1$ 是极小点。

再看
\[
f''(x)=6x.
\]
二阶导数在 $0$ 左负右正，故凹凸发生改变，$0$ 是拐点。注意 $f'(0)=-3\ne0$，拐点不要求驻点。

\subsection{同一个对象上的中值见证}

在任意 $a<b$ 上，令
\[
K=\frac{f(b)-f(a)}{b-a},
\qquad F(x)=f(x)-f(a)-K(x-a).
\]
则 $F(a)=F(b)=0$，所以存在 $\xi\in(a,b)$ 使 $f'(\xi)=K$。这说明“割线斜率被某条切线实现”并不依赖三次函数的特殊可解性，而是来自辅助函数与 Rolle 的通用结构。

\subsection{典型错误路径}

若只解 $f'(x)=0$ 并写“得到两个极值”，就漏掉了符号变化验证；若只解 $f''(x)=0$ 并写“得到拐点”，就漏掉了凹凸变化验证；若在中值证明中直接指定 $\xi=(a+b)/2$，就把存在性见证错误地固定成坐标。三处错误分别丢失了 Candidate/Confirm、区间符号和中值点存在性边界。

% ============================================================
\section{概念边界：真正判据、题目信号与错误后果}
% ============================================================

\begin{handbooklongtable}{L{0.13\linewidth}C{0.035\linewidth}L{0.13\linewidth}YY}
\toprule
\textbf{概念 A} & & \textbf{概念 B} & \textbf{为什么容易混 / 真正判据与题目信号} & \textbf{混淆后果}\\
\midrule
介值定理 & $\neq$ & Rolle & 介值只需连续并产生函数值；Rolle 还需内部可导和端点等值，产生导数零点。 & 用连续性直接推出内部水平切线\\
Rolle & $\neq$ & Lagrange & Lagrange 先减割线再用 Rolle；目标是平均变化率，不是只能证明零导数。 & 忘记构造割线修正\\
Lagrange & $\neq$ & Cauchy & Cauchy 同时运输两个函数的增量；比值形式还需分母资格。 & 未审查分母就相除\\
端点等值 & $\neq$ & 端点同号 & Rolle 需要 $F(a)=F(b)$；零点定理需要跨值/异号，二者用途不同。 & 把零点存在误当水平切线存在\\
零点 & $\neq$ & 重零点 & $F(a)=0$ 只计一个约束；还需 $F'(a)=0$ 等条件才可在反复 Rolle 中多次计数。 & 高阶导数见证少算或多算\\
驻点 & $\neq$ & 极值点 & $f'=0$ 只产生候选；一阶导数变号或充分条件才确认。 & 把平台拐点误判为极值\\
二阶导数为零 & $\neq$ & 拐点 & 拐点要求凹凸改变；$f''=0$ 只是候选，且不可导点也可能是拐点。 & 看到二阶零点就宣布拐点\\
严格割线序 & $\neq$ & 取极限后仍严格 & 严格不等式取极限一般只保留非严格次序；需先证 $f'$ 非减，再排除常值平台。 & 省略严格单调性的关键一步\\
局部导数符号 & $\neq$ & 全区间单调 & 单调结论需要对任意子区间调用 Lagrange 或其他全区间论证；单点符号没有全局信息。 & 从一个点的斜率外推整段函数\\
中值点存在 & $\neq$ & 中值点坐标已知 & 定理只给某个 $\xi$ 存在及其区间，通常不允许把它固定成对称点。 & 在证明中偷换见证点\\
直接积分 & $\neq$ & 已完成 Rolle 证明 & $F'=\Phi$ 只解决导数匹配；仍需端点等值或多零点。 & 以“找到原函数”替代存在性条件\\
\bottomrule
\end{handbooklongtable}

% ============================================================
\section{Topic04 的问题求解协议}
% ============================================================

面对中值定理证明、零点、单调极值或函数形状问题，按下列顺序维护状态：
\[
\boxed{
\begin{gathered}
\text{目标标准化}\to\text{资格审查}\to\text{辅助函数}\to\text{边界/零点设计}\\
\to\text{定理见证}\to\text{回译}\to\text{区间/符号验证}
\end{gathered}}
\]

\subsection{Step 1：Target}

把待证结论改写为 $\Phi(\xi)=0$、$f'(\xi)=K$、$F^{(n)}(\xi)=0$ 或函数值符号差。若问形状，先标出候选点与目标区间。

\subsection{Step 2：Qualification}

逐项检查闭区间连续、开区间可导、分母非零、对数真数、积分因子和端点是否属于定义域。资格失败时立即换定理或缩小合法区间。

\subsection{Step 3：Auxiliary Function}

逆用和积商链、割线修正、对数导数、积分因子或插值多项式，使某阶导数成为目标的非零倍数。辅助函数越短且端点越容易验证越好。

\subsection{Step 4：Boundary / Zero Design}

检查 $F(a)=F(b)$，或列出不同零点与重零点的区间顺序。需要多个中值点时分割区间；不能默认多次定理给出同一个点。

\subsection{Step 5：Witness and Translate Back}

明确引用介值、Rolle、Lagrange、Cauchy 或反复 Rolle，写出 $\xi$ 的所在区间，再展开 $F^{(k)}(\xi)=0$ 并除以已核对的非零因子。

\subsection{Step 6：Shape Verification}

做函数形状题时，按“候选集 → 一阶符号 → 二阶符号 → 端点比较”的顺序。独立检查包括：极值左右符号是否翻转、拐点凹凸是否改变、闭区间绝对极值是否比较端点、根的数量是否与 Rolle 下界相容。

% ============================================================
\section{高频失败路径：第一次偏离发生在哪里？}
% ============================================================

\begin{enumerate}
  \item \kw{资格失败}：未检查闭区间连续/开区间可导，就直接写“由中值定理”。
  \item \kw{目标失败}：没有先把结论标准化为导数或零点结构，盲目猜辅助函数。
  \item \kw{边界失败}：辅助函数导数匹配了目标，却没有制造端点等值或足够零点。
  \item \kw{分母失败}：Cauchy、商函数、对数导数或积分因子中分母可能为零，却直接相除。
  \item \kw{见证失败}：把存在的中值点固定成中点，或把多次应用得到的点误认为同一点。
  \item \kw{候选失败}：由 $f'=0$、$f''=0$ 或高阶导数为零直接宣布极值/拐点。
  \item \kw{严格性失败}：把严格不等式取极限后仍写成严格，未用平台反证恢复严格性。
  \item \kw{区间失败}：用一个点的导数符号替代整个区间的符号，或忘记闭区间端点比较。
  \item \kw{Owner 失败}：把插值余项、无限 Taylor、变限积分或完整 ODE 解法复制进本册。
\end{enumerate}

% ============================================================
\section{传统章节与 Source Corpus 映射}
% ============================================================

\begin{handbooklongtable}{L{0.16\linewidth}L{0.32\linewidth}Y}
\toprule
\textbf{Source / 传统内容} & \textbf{进入本册} & \textbf{分流与拒绝复制}\\
\midrule
I-07 定理链 & 介值/零点、Rolle、Lagrange、Cauchy、高阶 Rolle 的资格与生成关系 & 题目识别和首步动作进入 Rules\\
I-07 证明逆向工程 & 目标标准化、辅助函数反推、端点等值/零点设计与回译 & 只保留方法机制，不复制成题型清单\\
I-07 K 值法 & 割线修正、减去低次多项式、重零点计数 & 插值余项与误差估计移交 H-B02\\
I-07 乘积/商/对数/积分因子 & 作为辅助函数构造的局部结构接口 & 乘积/商本体归 Topic03；变限积分正则性归 Topic05/H-B03；ODE 解族归 Topic12\\
I-07 零点重数与反复 Rolle & 高阶导数见证、根数量下界和重根边界 & 多中值点插点法归 H-B02\\
I-06 函数形状 & 单调、极值、凹凸、拐点、零点重数与图像读法 & 曲率/局部二阶模型归 Topic03；定义域/函数对象归 Topic01\\
I-07.1 多中值点与中值点参数 & 只保留“分割区间保证不同中值点”的边界提示 & 主体归 H-B02，不在 Topic04 重写参数渐近\\
I-07.2 插值余项与导数估计 & 只保留“高阶 Rolle 接收插值条件”的接口 & Lagrange/Hermite 余项、Landau\textendash{}Kolmogorov 和误差界归 H-B02\\
\bottomrule
\end{handbooklongtable}

% ============================================================
\section{一页压缩：忘掉公式后怎样重建？}
% ============================================================

\begin{corebox}{一句话}
把目标表达式反推成辅助函数的导数，先验证区间资格，再制造端点等值或足够零点；中值定理只产生存在性见证，最后必须把见证回译并检查区间符号。
\end{corebox}

\subsection{母图}
\begin{processchain}
  \processstage{目标}
  \processstage{辅助函数}
  \processstage{资格}
  \processstage{端点/零点}
  \processstage{中值见证}
  \processstage{回译与形状}
\end{processchain}

\subsection{三个生成原则}
\begin{enumerate}
  \item \kw{运输原则}：介值运输函数值，Rolle 运输端点等值到内部零导数，Lagrange 运输割线斜率，Cauchy 运输双函数比值。
  \item \kw{设计原则}：辅助函数由目标导数结构与边界条件共同反推；只匹配导数而不制造边界约束是不完整的证明。
  \item \kw{候选原则}：驻点、二阶零点和中值点都是候选/存在性输出，必须用符号、区间和资格完成确认。
\end{enumerate}

\subsection{重建问题}
\begin{enumerate}
  \item 目标能否写成 $\Phi(\xi)=0$ 或某个导数等式？
  \item 需要哪个定理的资格：连续、可导、端点等值、非零分母还是多零点？
  \item 哪个辅助函数的导数最接近目标？它的端点值怎样验证？
  \item 得到的中值点位于哪个开区间，是否可能与另一个点混淆？
  \item $f'$ 或 $f''$ 的符号在区间上怎样变化？候选点是否真的确认？
  \item 若严格不等式经过了极限，严格性来自哪里？是否还需排除常值平台？
  \item 本题是否已经越界到 H-B02 的插值余项、Topic05 的积分接口或 Topic12 的 ODE 解族？
\end{enumerate}

\begin{center}\rule{0.5\linewidth}{0.5pt}\end{center}

\appendix
\section{中值定理与函数形状 Coverage 锚点}

\subsection{定理资格速查}

\begin{handbooktable}{L{0.16\linewidth}L{0.36\linewidth}Y}
\toprule
定理 & 资格 & 输出\\
\midrule
介值定理 & $[a,b]$ 连续 & 介于端点值之间的函数值存在\\
零点定理 & $[a,b]$ 连续且端点异号 & 内部零点存在\\
Rolle & 闭区间连续、开区间可导、端点等值 & $f'(\xi)=0$\\
Lagrange & 闭区间连续、开区间可导 & $f'(\xi)$ 等于割线斜率\\
Cauchy & 两函数均闭区间连续、开区间可导 & 交叉相乘的导数关系\\
高阶 Rolle & 足够阶光滑、零点/重零点条件 & 高阶导数零点存在\\
\bottomrule
\end{handbooktable}

\subsection{形状判定 Coverage}

按考纲查漏，本册覆盖：闭区间连续函数性质、介值与零点、Rolle\allowbreak/Lagrange\allowbreak/Cauchy 中值定理、单调性、极值、最大最小值、凹凸与拐点、函数图像和导数根的数量。Lagrange\allowbreak/Hermite 插值余项、导数估计和中值点参数渐近由 H-B02 负责；题面识别与路径选择进入 Rules。

\end{document}
